% ============================================================================
% rev4_section2_anatomy.tex
%
% Draft §2 "Anatomy of one AI facility gigawatt" for the rev-4 restructure.
% Standalone — not yet inserted into report.tex. Once approved, the rev-4
% restructure will renumber and reflow downstream sections.
%
% Every number resolves to a row in anatomy_layer_costs.yaml,
% anatomy_named_projects.yaml, or forecaster_capex_comparison.yaml.
% ============================================================================
\section{The anatomy of one AI facility gigawatt}
\label{sec:anatomy}

\textbf{Why a unit primitive.} Section~\ref{sec:scale} reports the headline aggregates: 7.37~GW tier-clean operational capacity, 49.8~GW raw 2030 horizon, 31.4~GW Monte~Carlo median, approximately \$1.84~trillion capital envelope, 78.0~million H100-equivalents. These quantities are scale numbers --- multiplications. The multiplicand is a single facility gigawatt of AI compute, and the answers to four small questions about that unit set the credibility of every aggregate downstream:

\begin{itemize}\itemsep2pt
\item What does one facility GW physically contain?
\item What does it cost to build, line item by line item?
\item What compute payload does it carry, and how does that change with each chip generation?
\item And against which denominator are the headline \$/GW numbers being quoted?
\end{itemize}

The last question is load-bearing. Apparent disagreements between forecasters of three to five times in \$/GW almost always disappear once the basis is reconciled. We dispatch it first.

\subsection{Reading every \$/GW figure: the four conventions}
\label{subsec:basis-conventions}

External forecasters quote AI data center capex in four distinct units. They are not interchangeable. Cross-source comparisons that ignore the basis distinction are arithmetic noise.

\begin{table}[h]
\centering
\small
\begin{tabular}{@{}L{4.6cm}R{1.7cm}R{1.7cm}L{6.0cm}@{}}
\toprule
\textbf{Convention} & \textbf{Low \$B/GW} & \textbf{High \$B/GW} & \textbf{Used by (representative)} \\
\midrule
Server-power GW, full all-in & 40 & 55 & Epoch AI headline (\$44B); NVIDIA / Jensen Huang (\$50--60B); Stargate (\$50B) \\
\textbf{Facility-power GW, full all-in} & \textbf{30} & \textbf{47} & \textbf{Epoch facility (\$29--30B); McKinsey (\$33B); Bernstein (\$28B PUE-adj.); \emph{this paper, central \$37B}} \\
IT-load GW, shell + M\&E only (no IT) & 10 & 15 & Cushman \& Wakefield (\$11.7M/MW); JLL (\$11.3M/MW); Meta Hyperion JV (\$13.5B/GW); Crusoe Abilene (\$12.5M/MW) \\
AI fit-out delta (incremental over standard hyperscale) & +20 & +30 & JLL (+\$25M/MW); ConstructElements (+\$20M/MW) \\
\bottomrule
\end{tabular}
\caption{The four conventions in which AI data center capex is quoted in 2025--2026. The first two differ only in denominator (server-power MW vs.\ facility-power MW, related by the PUE ratio of roughly 1.10--1.30). The last two exclude either the IT hardware or the entire base hyperscale envelope. \emph{This paper uses the second convention throughout}: per facility-power GW, with IT hardware included and behind-the-meter generation reported separately. All eighteen sources canvassed and their decompositions appear in \texttt{forecaster\_capex\_comparison.yaml}.\protect\footnotemark}
\label{tab:basis-conventions}
\end{table}
\footnotetext{Sources canvassed: Epoch AI \emph{Frontier Data Centers Hub} (Nov.~2025); Bernstein (Stacy Rasgon, Nov.~2025, via Investing.com); McKinsey \emph{Cost of Compute} (Apr.~2025); Bain \emph{Compute Power Report} (Sep.~2025); JLL \emph{2026 Global Data Center Outlook}; Cushman \& Wakefield \emph{2025 Cost Guide}; Goldman Sachs (Sheridan/Davenport, 2025); Morgan Stanley (Stephen Byrd, 2025); JP Morgan; BloombergNEF; Synergy Research; Dell'Oro Group; IEA \emph{Energy and AI}; LBNL \emph{2024 US Data Center Energy Usage Report}; Aschenbrenner \emph{Trillion-Dollar Cluster}; NVIDIA Q2 FY26 earnings; Barclays; OpenAI/Oracle Stargate; Meta--Blue Owl Hyperion JV; CoreWeave 10-K FY25; Nebius Q4 FY25. Full URLs and retrieval dates in \texttt{forecaster\_capex\_comparison.yaml}.}

The basis confusion is not academic. Epoch AI's headline \$44~B/GW figure (server-power basis) and McKinsey's \$33~B/GW figure (facility-power, incl.\ infrastructure) differ by roughly the PUE ratio --- they are saying close to the same thing in different units.\footnote{Epoch AI, \emph{Frontier Data Centers Hub}, retrieved 2026-04-27, \url{https://epoch.ai/blog/introducing-the-frontier-data-centers-hub/}. The Epoch documentation breaks the \$44B/GW into \$30B IT hardware plus \$14B other; the facility-GW figure is approximately \$29--30B at PUE~1.20. McKinsey \& Company, \emph{The Cost of Compute: A \$7~Trillion Race to Scale Data Centers}, April~2025, \url{https://www.mckinsey.com/industries/technology-media-and-telecommunications/our-insights/the-cost-of-compute-a-7-trillion-dollar-race-to-scale-data-centers}.} NVIDIA's \$50--60~B/GW ``AI factory'' figure is server-power basis and self-interested but inside the same band. The Crusoe--Oracle--OpenAI Stargate Abilene disclosure of \$15~B for 1.2~GW (\$12.5M/MW) is correctly an order of magnitude smaller because it excludes the IT hardware --- Oracle funds the GPU layer separately, on its own balance sheet, at an additional reported \$40~B for approximately 400{,}000 GB200s.\footnote{Crusoe Energy Systems, ``Crusoe expands AI data center campus in Abilene to 1.2 gigawatts'', 2025-05, \url{https://www.crusoe.ai/resources/newsroom/crusoe-expands-ai-data-center-campus-in-abilene-to-1-2-gigawatts}; Bloomberg, ``JPMorgan leads \$7.1~billion loan for Blue Owl-tied data center'', 2025-05-22; Data Center Dynamics, ``Crusoe secures \$11.6bn in debt and equity for OpenAI's Stargate data center campus'', 2025-05. The Oracle GPU figure (\$40B for ~400K GB200s) is reported via Spyglass/SemiAnalysis and not Crusoe-confirmed.}

\subsection{The physical bill of materials}
\label{subsec:physical-bom}

A 2026-vintage 1~GW AI campus is not a building --- it is a power plant with computers attached. Six physical layers, each with a distinct supply chain, lead time, and inflation profile, share the construction-side ledger; the accelerator and server bill of materials sits on top.

\begin{table}[h]
\centering
\small
\begin{tabular}{@{}L{4.6cm}R{1.5cm}R{1.5cm}R{1.6cm}L{4.4cm}@{}}
\toprule
\textbf{Layer} & \textbf{Low \$M/MW} & \textbf{High \$M/MW} & \textbf{Central \$M/MW} & \textbf{Layer-dominating components} \\
\midrule
Shell + civil + land + base/AI fit-out & 4.4 & 9.7 & 6.5 & Building shell; AI fit-out delta (CDU plumbing, busway) \\
Cooling stack & 2.0 & 3.5 & 2.75 & DLC cold plates and CDUs; chillers; dry coolers in water-stressed sites \\
Power infrastructure (grid-tied) & 1.9 & 3.7 & 2.8 & Main transformers (long lead); UPS \\
\quad\emph{or with behind-the-meter gas} & \emph{2.5} & \emph{5.5} & \emph{4.0} & \emph{Above plus 1.0--2.5 \$M/MW gas turbines} \\
Networking fabric (rack-internal scale-up) & --- & --- & \emph{(in rack BOM)} & NVLink scale-up; copper cartridges \\
Out-of-rack networking & 2.5 & 4.5 & 3.5 & Optics (60\% of out-of-rack BOM); back-end fabric \\
Accelerator + server BOM (rack-complete) & 16 & 23 & 19.5 & Accelerator silicon; HBM \\
Storage & 0.3 & 0.8 & 0.55 & NVMe arrays; controllers \\
Grid interface (substation-to-fence + assigned upgrades) & 0.3 & 5.0 & 1.1 & Wide range by ISO posture; see \S\ref{subsec:scaling-the-unit} \\
\midrule
\textbf{Subtotal: facility-only (no IT)} & \textbf{11} & \textbf{19} & \textbf{13.2} & Reconciles to Crusoe Abilene \$12.5M/MW disclosed \\
\textbf{Subtotal: IT BOM at customer} & \textbf{19} & \textbf{28} & \textbf{23.5} & Accel + server + out-of-rack networking + storage \\
\textbf{Total all-in facility-GW (2026)} & \textbf{30} & \textbf{47} & \textbf{37} & Bernstein \$35B central; JLL \$45--55B fully built \\
\bottomrule
\end{tabular}
\caption{Cost decomposition for one 2026-vintage AI facility gigawatt, by physical layer. Central case is grid-tied with diesel backup; the BTM-gas alternative adds \$1.0--2.5M/MW. Every sub-component breakdown, with primary citations and evidence tier (T1--T6) per row, lives in \texttt{anatomy\_layer\_costs.yaml}.}
\label{tab:layer-bom}
\end{table}

\textbf{Where the cost moves the most}: the cooling layer ratio rose from a historical 15--20\% of facility capex to 20--30\% for 2026 builds, almost entirely driven by the GB200/GB300 transition from 30--50~kW air-cooled racks to 80--132~kW direct-liquid-cooled racks.\footnote{Direct liquid cooling (DLC) revenue surged 156\% year-over-year in 2Q~2025 (Dell'Oro Group); the total liquid cooling market exceeded \$2.5~B in 2025 and is forecast above \$8~B by 2030. NVIDIA's NVL72 reference architecture mandates 25--45\textdegree C inlet liquid, 80~L/min flow, junction temp $<$75\textdegree C --- physically incapable of being held by air at 1.2~kW per GPU. Schneider Electric's \$850M acquisition of Motivair (October~2024) and AWS's launch of in-house IRHX direct-to-chip cooling (July~2025) are the two most significant supply-chain moves. Sources: Vertiv press release on the NVL72 7~MW reference architecture (October~2024); Dell'Oro DCPI Q4 2025; ASHRAE TC9.9 liquid cooling bulletin (September~2024).} Power transformers --- a single line item of \$300--600k per MW --- have \emph{prices up 77\% since 2019} and \emph{lead times of 80--128 weeks for power transformers and 144 weeks for generator step-up units} (Wood Mackenzie 2025).\footnote{Wood Mackenzie via PowerMag, ``Transformers in 2026: shortage, scramble, or self-inflicted crisis?'', August~2025, \url{https://www.powermag.com/transformers-in-2026-shortage-scramble-or-self-inflicted-crisis/}. NERC 2025 Summer Reliability Assessment confirms 2024 average lead times of approximately 120 weeks with large-class units at 80--210 weeks.} Behind-the-meter gas, when installed, doubles the power-infrastructure layer cost on its own --- \$1.0--2.5M/MW capex plus the supply-chain reality that GE Vernova ended 2025 with an 80~GW gas-turbine backlog stretching into 2029, with reservation slots out to 2031.\footnote{GE Vernova Q3 2025 earnings call, October~22, 2025; Utility Dive, ``GE Vernova's gas turbine backlog stretches into 2029'', September~2025, \url{https://www.utilitydive.com/news/ge-vernova-gas-turbine-investor/807662/}. Siemens Energy reports a similar backlog into 2030 in its FY2025 annual report.}

\textbf{Worked example: Stargate Abilene.} Crusoe Energy's disclosed all-in capital for the 1.2~GW Abilene campus is \$15~billion (\$12.5M/MW), composed of \$9.6~B JPMorgan-led senior debt (a \$2.3~B initial facility plus a \$7.1~B construction loan announced May~2025) and approximately \$5~B equity from Crusoe, Blue~Owl Real~Assets, and Primary~Digital~Infrastructure. This figure represents shell + power + cooling + site infrastructure only. Oracle, the offtake tenant on a 15-year lease, funds the IT layer separately at a reported additional \$40~B for approximately 400{,}000 GB200 GPUs. The campus runs closed-loop direct-to-chip liquid cooling with approximately one million gallons per building of initial water fill and approximately 50{,}000 gallons per building per year for water-quality maintenance --- effectively zero evaporative consumption --- with twenty-nine GE Vernova LM2500XPRESS aeroderivative gas turbines (\textasciitilde35~MW each, \textasciitilde1~GW combined) providing behind-the-meter generation that supplements ERCOT power. Phase~1 was delivered twelve months from groundbreaking, against an industry-standard bid of two-and-a-half years.\footnote{Crusoe Energy Systems newsroom releases (2024--2026); Madrona ``Founded \& Funded'' interview with Chase Lochmiller (CEO), retrieved 2026-04-27, \url{https://www.madrona.com/the-infrastructure-of-intelligence-inside-crusoes-ai-factory-in-texas/}. The 15-year Oracle lease and OpenAI sub-tenancy structure are confirmed by Lochmiller in DCD: ``Our customer is Oracle. OpenAI is Oracle's customer.'' Annual rent of approximately \$1~billion is reported by SemiAnalysis (secondary; not Crusoe-confirmed).}

The Crusoe disclosure is the cleanest landlord/tenant capex split publicly available and \emph{independently reconciles} to the \$13.2M/MW central case for the facility-only-no-IT subtotal in Table~\ref{tab:layer-bom}. Behind-the-meter gas at Abilene is amortized through the lease structure rather than landed in shell capex, which is why Crusoe's disclosed number sits slightly below the bottoms-up table's BTM-inclusive central case of \$15.4M/MW.

\subsection{The compute payload: chips per gigawatt is falling, compute per gigawatt is rising}
\label{subsec:compute-payload}

The single most counterintuitive finding from the bottoms-up reconstruction is that \textbf{the number of chips per facility-MW \emph{decreases} with each chip generation}. Power per chip is rising faster than performance per watt at the rack level: an H100 SXM is 700~W, a GB200 superchip is 2{,}700~W including the Grace partner, and a Rubin Ultra die is rumored above 1.5~kW.\footnote{NVIDIA \emph{H100 Tensor Core GPU} datasheet (700~W TDP); Tom's Hardware on GB200 NVL72 architecture (2025-03-21); NVIDIA GTC 2026 keynote on Vera Rubin and Rubin Ultra (Kyber NVL576) generation. The 600~kW Kyber rack target was disclosed at GTC March~2026.}

\begin{table}[h]
\centering
\small
\begin{tabular}{@{}L{4.0cm}R{2.4cm}R{1.6cm}R{2.0cm}R{2.0cm}@{}}
\toprule
\textbf{Vintage / platform} & \textbf{Power per chip (system-allocated)} & \textbf{Chips / facility-MW} & \textbf{H100e per chip} & \textbf{H100e / facility-MW} \\
\midrule
2024 — HGX H100 (8-GPU air-cooled) & 1{,}389~W & \textasciitilde570 & 1.00 & \textasciitilde0.57M \\
2025--26 — GB200 NVL72 (liquid) & \textasciitilde1{,}900~W & \textasciitilde440 & 2.50 & \textasciitilde1.10M \\
2026 — GB300 NVL72 & \textasciitilde2{,}100~W & \textasciitilde415 & 3.75 & \textasciitilde1.55M \\
2H'2026--2027 — Vera Rubin NVL144 & \textasciitilde1{,}050~W & \textasciitilde850 & 7.50 & 6.4M (theoretical pure) \\
2028 (realistic deployed mix) & --- & 700--800 & 8--10 & 1.5--2.5M (installed-base wgt.) \\
\bottomrule
\end{tabular}
\caption{Chip and compute density per facility-MW by silicon vintage, at PUE~1.15 with 10\% network/storage overhead. The 2028 row reflects an installed-base-weighted mix (Blackwell-Ultra installed base with a Rubin tail), which is the appropriate forecasting unit; pure-Rubin theoretical density is roughly 4$\times$ higher but will not characterize deployed fleets before 2029. SemiAnalysis's reported 100k-H100 cluster benchmark of 1{,}389~W per H100 system-allocated (20{,}480 GPUs in 28.4~MW critical-IT) reconciles to approximately 720 H100/IT-MW and 600 H100/facility-MW. Per-chip H100e ratios from NVIDIA-published peak FP8 throughput.\protect\footnotemark}
\label{tab:chips-per-mw}
\end{table}
\footnotetext{SemiAnalysis ``100,000 H100 Clusters: power, network topology, ethernet vs.\ infiniband, reliability'' (cluster benchmark 28.4~MW for 20{,}480 GPUs); NVIDIA datasheets for H100, B200, GB200, GB300; NVIDIA GTC 2026 disclosures for Rubin and Rubin Ultra. Full per-chip table with 18 accelerators (NVIDIA, AMD MI300X--MI400, Google TPU v5p--Ironwood, AWS Trainium 2--3, Microsoft Maia 100--200, Meta MTIA v2--v4) lives in \texttt{anatomy\_layer\_costs.yaml}, layer \texttt{accelerator\_server\_bom}.}

The compute story is therefore not ``more chips per watt'', it is ``more H100-equivalents per chip''. Holding the buildout constant, the H100e-per-facility-MW figure roughly triples between 2024 and 2026 and rises by another 50\% to 2028. Section~\ref{sec:silicon} traces this through the chip roadmap and validates the per-vintage densities against deployed clusters.

\textbf{The accelerator capex line.} The bottoms-up reconstruction --- approximately 440 GB200-class chips per MW at \textasciitilde\$45k all-in chip-plus-rack-share --- yields approximately \$19--22~B/GW for the accelerator and server bill of materials on a 2026 all-NVIDIA build. A hyperscaler-mix build with 40--50\% custom-silicon share (Trainium~3 at approximately \$0.65 per H100e of FP8 throughput, TPU Ironwood at approximately \$1.20 per H100e, Maia~200 at approximately \$1.30 per H100e, all compared to GB200's approximately \$2.50 per H100e) lands in the \$15--18~B/GW range. The blended Western 2026 horizon, weighting the operator mix (Microsoft / Amazon / Google / Meta custom-silicon-heavy at 40--50\%; OpenAI / xAI / neoclouds at all-NVIDIA), sits at approximately \textbf{\$18--21~B/GW for the rack-complete accelerator+server BOM}, with the headline central case of \$19.5~M per facility-MW.

\subsection{The financial bill of materials}
\label{subsec:financial-bom}

Aggregating Table~\ref{tab:layer-bom} for a typical 2026-vintage Tier-2/3 build (Texas, Ohio, Mississippi, Iowa, Wisconsin) gives a \textbf{central all-in facility-GW capex of approximately \$37~billion}, with a defensible band of \$30--47~B/GW spanning the realistic configuration space. This number sits at the center of the external literature distribution.

\begin{table}[h]
\centering
\small
\begin{tabular}{@{}L{5.4cm}R{2.4cm}L{6.0cm}@{}}
\toprule
\textbf{External anchor} & \textbf{\$B/facility-GW} & \textbf{Reconciliation} \\
\midrule
Crusoe Abilene (shell + power + cooling, no IT) & 12.5 & Matches our facility-only-no-IT central (\$13.2M/MW) within 5\% \\
Meta Hyperion / Blue Owl JV (shell + M\&E only) & 13.5 & Matches our facility-only-no-IT central directly; IT and Entergy gas separate \\
Bernstein (Stacy Rasgon, IT-load basis) & 35 & PUE-adjusted to \$28B/facility-GW; below our central; under-counts BTM gas \\
Epoch AI facility-GW & 29--30 & Below our central; consistent with Epoch using lower-density chip mix \\
McKinsey \emph{Cost of Compute} (derived) & 33 & Within rounding distance of our central \\
JLL fully-built ecosystem (high) & 45--55 & Aligns with our high-end at \$47B \\
NVIDIA / Jensen ``AI factory'' (server-power basis) & 50--60 & Translates to \$40--48B/facility-GW; self-interested high \\
\bottomrule
\end{tabular}
\caption{Reconciliation of the bottoms-up \$30--47~B/facility-GW range to seven external anchors. The Crusoe Abilene and Meta Hyperion JV disclosures are the cleanest direct anchors because both are \emph{actual closed financings} and both explicitly exclude the IT layer, which makes them directly comparable to the facility-only-no-IT subtotal in Table~\ref{tab:layer-bom}. All anchor URLs in \texttt{forecaster\_capex\_comparison.yaml}.}
\label{tab:capex-reconciliation}
\end{table}

\subsection{Scaling the unit}
\label{subsec:scaling-the-unit}

Multiplying one facility GW by the announced horizon resolves the headline numbers. The Western 2026--2030 capital envelope at the raw 51.9~GW announced horizon is approximately \textbf{\$1.9~trillion central, with a defensible range of \$1.6--2.4~trillion} depending on operator mix, BTM-generation share, and silicon-vintage assumptions.

\begin{table}[h]
\centering
\small
\begin{tabular}{@{}L{6.4cm}R{2.0cm}R{3.6cm}@{}}
\toprule
\textbf{Horizon point} & \textbf{Facility GW} & \textbf{Capex-equivalent at \$30--47B/GW} \\
\midrule
Monte~Carlo p10 (\S\ref{subsec:monte-carlo}) & 24.6 & \$0.74--1.16T \\
Scenario-adjusted median (Monte~Carlo p50) & 31.8 & \$0.95--1.49T \\
Deterministic tier-weighted point & 36.8 & \$1.10--1.73T \\
Monte~Carlo p90 & 37.4 & \$1.12--1.76T \\
Raw non-stretch horizon (announced minus T5) & 45.2 & \$1.36--2.12T \\
\textbf{Raw announced horizon} & \textbf{51.9} & \textbf{\$1.56--2.44T (central \$1.9T)} \\
Full-realization ceiling & 55.3 & \$1.66--2.60T \\
\bottomrule
\end{tabular}
\caption{Capex-equivalent at the bottoms-up \$30--47~B/facility-GW range across the seven horizon points used elsewhere in this paper. The figures are \emph{capex commitments implied by the announced build}, not \emph{cash energized by 2030}. A campus that slips its commercial operating date may still consume equipment, shell, and grid capex before 2030 even if it is not energized; conversely, the \$30--47~B/GW figure assumes 2026-vintage component prices and a Tier-2/3 build that may not generalize to extreme cases (Northern Virginia land, behind-the-meter nuclear PPAs, or pure-Rubin builds).}
\label{tab:scaling-the-unit}
\end{table}

\textbf{Capex-equivalent versus realized energized capacity.} The \textbf{\$1.9~trillion (range \$1.6--2.4~T) capital envelope} is the capex commitment implied by the raw 51.9~GW announced horizon at the central \$/GW. The Monte~Carlo median realized horizon of 31.8~GW implies, at the same per-GW range, a smaller realized energized-capacity capital base of approximately \$1.0--1.5~T --- but this is not necessarily a smaller cash spend, because campuses that slip their commercial operating dates may still absorb equipment, shell, and grid capex before 2030. The two numbers measure different things: announced build envelope versus realized energized capacity. The capital envelope is therefore best read as a commitment surface, not an energized-by-2030 fleet, and is one of several places in this paper where headline aggregates and probability-weighted realizations diverge.

\textbf{What this section anchors for the rest of the paper.} Section~\ref{sec:scale} reports the GW aggregates; Section~\ref{sec:geography} asks where one such GW can be physically delivered; Section~\ref{sec:operators} asks who controls them; Section~\ref{sec:silicon} asks what compute payload each one carries; Section~\ref{sec:capital} asks who funds the multiplication and who underwrites the revenue (the \$1.9~T total against \$550~B of contracted RPO); Section~\ref{sec:execution-risk} asks what prevents a gigawatt from becoming operational. Each of those sections is now multiplying or interrogating a unit that has been priced and assembled here.
